%% ma: L12-B01-0046
%% lop: 12
%% chuong: 1
%% bai: 1
%% dang: 12.1.16
%% loai: TN4
%% muc: TH
%% dap_an: "A"
%% nguon: "Ảnh giáo viên gửi 05/11/2026 (M11): Kiểm tra giữa HK1 Toán 12 (2026-2027), THPT Hồ Thị Bi (TP.HCM), mã đề 1133, Phần I Câu 22"
%% kien_thuc: "Bài 1 (tính đơn điệu của hàm số phân thức bậc nhất trên bậc nhất)"
%% trang_thai: chinh-thuc
%% ghi_chu: "Đề gốc không kèm đáp án; đã giải độc lập."
\begin{ex}
Cho hàm số $y=\dfrac{2x-1}{1-x}$. Khẳng định nào dưới đây đúng?
\choice{\True Hàm số đồng biến trên các khoảng $(-\infty;1)$ và $(1;+\infty)$}{Hàm số nghịch biến trên khoảng $(-1;1)$}{Hàm số đồng biến trên $\mathbb R$}{Hàm số nghịch biến trên các khoảng $(-\infty;-1)$ và $(-1;+\infty)$}
\loigiai{Tập xác định $\mathbb R\setminus\{1\}$. $y'=\dfrac{2(1-x)+(2x-1)}{(1-x)^2}=\dfrac1{(1-x)^2}>0$ với mọi $x\ne1$.
Vậy hàm số đồng biến trên $(-\infty;1)$ và $(1;+\infty)$.
Hàm số không xác định tại $x=1$ nên không đồng biến trên $\mathbb R$.}
\end{ex}
